\scanpage{266}


\diagram{\makebox[0pt][r]{$(148)$\quad}\begin{tikzcd}[ampersand replacement=\&,column sep=large]
\widetilde Z \& Z \arrow[l,"c'"'] \arrow[d,"u\,{\unsure{?|\sim|\wr}}"']\\
\deleted{T_0}\arrow[u,"\deleted{u_0}"'] \& T_H\arrow[ul,"\unsure{c'_0}"]\arrow[l,"\deleted{\unsure{c'_0}}"]
\end{tikzcd}}

\[
\operatorname{int}_H(t^{-1})=c'_0(t^{-1})\longleftarrow\joinrel\shortmid\ t
\]

i.e. qu'on doit avoir

\[
\operatorname{int}_H(u(\lambda))^{-1}\deleted{\unsure{\vert_H}\,\unsure{\operatorname{int}}}
=\operatorname{int}\bigl(u_0(c'(\lambda))\bigr)|H
\]

ou encore

\[
u(\lambda)^{-1}\equiv u_0(c'(\lambda))\quad\bmod\ \operatorname{Cent}_G(H)
\]

\deleted{\unsure{ou on a|i.e. que}}\inserted{i.e. que}

\[
u(\lambda).u_0(c'(\lambda))\in\operatorname{Cent}_G(H)
\]

or on a (128)

\[
u(\lambda)\,u_0(c'(\lambda))=\lambda\unsure{\in|+}Z\in\operatorname{Cent}(G),
\]

ce qui \deleted{\unsure{vient}} prouve (148).

Cela achève donc d'expliciter ce qu'il fallait via $T_0$. On aurait l'explicitation symétrique via $T_1$, avec $U_1$ remplaçant $U_0$ -- mais \inserted{en fait on trouve \uline{le même} sous-groupe de $\operatorname{Aut}(H)$ en fait} \deleted{en passant à $T_G$, on fera attention que $T_G$ n'opère \uline{pas} \uline{fidèlement} sur $H$, donc $T_G$ ne peut s'expliciter comme sous-groupe de $\operatorname{Aut} H$.}

\[
\begin{aligned}
\widetilde Z &=\operatorname{Im}(T_0\to\operatorname{Aut}(H))=\operatorname{Im}(T_1\to\operatorname{Aut}(H))\deleted{\unsure{=\operatorname{Im}(T)}}\\
&=\operatorname{Im}(T_G\to\operatorname{Aut} H)
\end{aligned}\tag*{(149)}
\]

\deleted{Mais} Explicitons encore le sous-groupe $\unsure{\mathfrak z|\mathcal Z}\inserted{=\operatorname{Ker}(\delta_{T_0}:T_0\to Z)}$ de $T_0$ en tant que sous-groupe de $\operatorname{Aut} H$, on trouve \unsure{p}

\[
\begin{aligned}
\mathfrak z&\xrightarrow{\sim}\quad\text{sous-groupe de }\operatorname{Aut}H\ \text{formé des}\\
z&\longmapsto\operatorname{int}(z)|H\quad \text{autom }\bar g\ \text{qui normalisent }U_0,\\
&\hspace{8em}\text{et qui centralisent }T_H\ \text{et }\sigma.\\
&\hspace{8em}\deleted{\text{(donc aussi }\sigma(U_0)\text{)}}\ \text{qui}
\end{aligned}\tag*{(149)}
\]
% archive sheet number: 265
